\chapter{Kết luận và hướng phát triển}
\label{chap:conclusion}

\section{Kết luận}
Đề tài đã xây dựng một hệ thống đa tác tử LLM cho tín hiệu giao dịch Bitcoin, trong đó mỗi nguồn dữ liệu do một agent
chuyên trách phân tích, mâu thuẫn giữa các agent được đo bằng chỉ số kết hợp phân kỳ KL và phương sai, và khi mâu thuẫn
cao, hệ thống phân tích nguyên nhân gốc và tổ chức tranh biện có kiểm tra nguồn trích dẫn. Hệ thống được đánh giá đối
chứng với CryptoTrade, một baseline đã công bố tại EMNLP 2024, được cài đặt lại và kiểm chứng khớp với mã nguồn gốc,
trên cùng dữ liệu, cùng môi trường giao dịch và cùng mô hình ngôn ngữ chạy cục bộ.

Câu trả lời cho ba câu hỏi nghiên cứu:
\begin{description}[leftmargin=1.2cm,style=nextline]
  \item[RQ1] Ở dạng hiện tại, hệ thống đề xuất (+0,38\% trung bình) chưa vượt CryptoTrade đầy đủ (+5,71\%), dù vượt
  phiên bản CryptoTrade chỉ dùng dữ liệu thị trường ($-$3,35\%). Tỷ lệ dự đoán đúng hướng tương đương CryptoTrade; chênh
  lệch lợi suất chủ yếu đến từ mức vào lệnh cố định.
  \item[RQ2] Giá trị của từng nguồn đổi dấu theo giai đoạn thị trường: tin tức hữu ích khi giá tăng, on-chain hữu ích
  khi đi ngang, và cả hai làm kết quả xấu đi khi giá giảm.
  \item[RQ3] Tranh biện giảm mức mâu thuẫn từ khoảng 0,61 xuống 0,22 và không làm kết quả xấu đi, nhưng chỉ đổi 11
  trên 195 quyết định và chỉ theo hướng đứng ngoài thị trường, do thiết kế không cho phép agent đổi hướng.
\end{description}

Các kết quả này cho thấy việc xử lý mâu thuẫn tường minh có tác dụng giảm rủi ro và tăng khả năng giải thích (mỗi lập
luận đều có nguồn kiểm chứng được), nhưng để cải thiện lợi suất, hệ thống cần biết \emph{khi nào nên tin nguồn nào}
thay vì dùng trọng số cố định.

\section{Hạn chế}
\begin{itemize}
  \item Đánh giá trên một tài sản, một mô hình ngôn ngữ và ba giai đoạn ngắn; chưa có kiểm định ý nghĩa thống kê.
  \item Các tham số của Conflict Analyzer, tranh biện và Mediator chưa được hiệu chỉnh trên dữ liệu.
  \item Hệ số phạt thời gian trong Mediator vô tình làm giảm khoảng 60\% trọng số của on-chain và tin tức so với dữ liệu
  giá (Mục~\ref{sec:mediator}).
  \item Thời gian chạy của hệ thống đề xuất gấp khoảng 5 lần CryptoTrade đầy đủ.
\end{itemize}

\section{Hướng phát triển}
\paragraph{Trong ngắn hạn (đang thực hiện).} Xây dựng một bộ nhận diện xu hướng chỉ dựa trên dữ liệu giá tới ngày $D-1$
(ví dụ giá so với MA50, lợi suất 20--30 ngày) và dùng cho hai mục đích:
\begin{enumerate}
  \item \textbf{Trọng số agent theo xu hướng}: Mediator điều chỉnh mức ảnh hưởng của từng agent theo giai đoạn tăng,
  giảm hoặc đi ngang, khai thác trực tiếp phát hiện của RQ2. Hướng này gần với Market Regime Council~\cite{mrc2026};
  điểm khác là đề tài kết hợp trọng số theo xu hướng với cơ chế phát hiện và giải thích mâu thuẫn.
  \item \textbf{Bộ lọc xu hướng}: chặn tín hiệu đi ngược xu hướng rõ ràng.
\end{enumerate}
Để tránh khớp quá mức với tập kiểm thử, tham số sẽ được chọn trên khoảng 150 ngày dữ liệu CryptoTrade chưa sử dụng, sau
đó mới đánh giá trên ba giai đoạn kiểm thử, và kết quả của mọi biến thể đều được báo cáo. Cả hai hướng có thể thử
nghiệm trên các nhận định đã lưu mà không cần gọi lại mô hình.

\paragraph{Trong dài hạn.}
\begin{itemize}
  \item Xác định mức vào lệnh theo độ tự tin hoặc theo $\Sfinal$ thay vì cố định 50\%.
  \item Cho phép agent đổi hướng sau tranh biện khi lập luận phản biện đủ mạnh và có nguồn kiểm chứng, để tranh biện trở
  thành cơ chế sửa sai chứ không chỉ giảm rủi ro.
  \item Mở rộng đánh giá sang ETH, SOL như bài báo gốc, thêm mô hình ngôn ngữ khác và kiểm định ý nghĩa thống kê.
\end{itemize}
